%% ============================================================
%%  SECCIÓN 3 — Objetivos
%% ============================================================

\section{Objetivos}
\label{sec:objetivos}

\subsection{Objetivo General}
Familiarizarse con los principios de operación, calibración y manejo del osciloscopio digital como instrumento de medición de diferencias de potencial continuo (DC), tensiones alternas variables en el tiempo (AC), parámetros fundamentales de formas de onda periódicas (amplitud, período y frecuencia) y como sistema graficador paramétrico bidimensional $X$-$Y$.

\subsection{Objetivos Específicos}
\begin{itemize}
    \item Identificar y operar los controles principales del panel frontal del osciloscopio digital (acoplamientos $\text{GND}$, $\text{DC}$ y $\text{AC}$, sensibilidad vertical en \unit{\volt\per\division}, base de tiempo horizontal en \unit{\second\per\division} y controles de disparo o \textit{trigger}).
    \item Medir tensiones continuas provenientes de pilas electroquímicas y fuentes reguladas multisalida a distintas escalas verticales, evaluando el criterio de escala óptima para la reducción del error de apreciación.
    \item Comparar cuantitativamente las lecturas obtenidas mediante la retícula del osciloscopio con los valores patrones registrados por un multímetro digital True-RMS de alta precisión.
    \item Determinar experimentalmente el voltaje pico ($V_p$), voltaje pico a pico ($V_{pp}$), período ($T$) y frecuencia ($f$) de la señal senoidal entregada por un transformador reductor de \qty{220}{\volt} a \qty{6}{\volt} a \qty{60}{\hertz}.
    \item Caracterizar señales alternas periódicas sintéticas (senoidales, triangulares y cuadradas) generadas por un generador de funciones arbitrarias a \qty{250}{\hertz}, verificando la relación entre amplitud pico y voltaje eficaz.
    \item Configurar el osciloscopio en modo graficador $X$-$Y$ para observar y analizar figuras de Lissajous para relaciones de frecuencia $1:1$, $1:2$, $1:3$ y $1:4$, así como el comportamiento de señales constantes DC bajo inversión polar del canal vertical.
\end{itemize}
